\begin{table}[H]\centering
\caption{Total amount billed per delivery, by mode}
\label{tab:cost}
\small
\small\setlength{\tabcolsep}{4pt}\resizebox{\ifdim\width>\linewidth \linewidth\else\width\fi}{!}{%
\begin{tabular}{lcc}
\toprule
 & Cesarean & Vaginal \\
\midrule
Mean billed cost (R\$) & 6,739 & 6,743 \\
Median billed cost (R\$) & 5,417 & 5,178 \\
Deliveries, 2015--2024 & 2,520,552 & 480,604 \\
\bottomrule
\end{tabular}
}
\\[2pt]\footnotesize\textit{Notes:} TISS delivery hospitalizations, 2015--2024. Total billed cost sums all billed items (procedures, materials, drugs, daily rates) under each delivery event; it is the charged/informed value, not the negotiated price paid, so it is a gross order-of-magnitude figure. The top and bottom 0.5\% of the cost distribution are trimmed. A cesarean and a vaginal delivery bill nearly the same total amount, so the epidemic is not explained by higher cesarean billing.
\end{table}
